\documentclass[10pt,a4paper]{amsart}
\usepackage{amsmath,amssymb,amsthm,graphicx}
\usepackage{verbatim}
\newenvironment{reference}{\comment}{\endcomment}
%\newenvironment{reference}{}{}
\newenvironment{slogan}{\comment}{\endcomment}
\newenvironment{history}{\comment}{\endcomment}

% For commutative diagrams we use Xy-pic
\usepackage[all]{xy}

% We use 2cell for 2-commutative diagrams.
\xyoption{2cell}
% Theorem environments.
%
\theoremstyle{plain}
\newtheorem{theorem}[subsection]{Theorem}
\newtheorem{proposition}[subsection]{Proposition}
\newtheorem{lemma}[subsection]{Lemma}

\theoremstyle{definition}
\newtheorem{definition}[subsection]{Definition}
\newtheorem{example}[subsection]{Example}
\newtheorem{exercise}[subsection]{Exercise}
\newtheorem{situation}[subsection]{Situation}

\theoremstyle{remark}
\newtheorem{remark}[subsection]{Remark}
\newtheorem{remarks}[subsection]{Remarks}

\numberwithin{equation}{subsection}

% Macros
%
\def\lim{\mathop{\mathrm{lim}}\nolimits}
\def\colim{\mathop{\mathrm{colim}}\nolimits}
\def\Spec{\mathop{\mathrm{Spec}}}
\def\Hom{\mathop{\mathrm{Hom}}\nolimits}
\def\Ext{\mathop{\mathrm{Ext}}\nolimits}
\def\SheafHom{\mathop{\mathcal{H}\!\mathit{om}}\nolimits}
\def\SheafExt{\mathop{\mathcal{E}\!\mathit{xt}}\nolimits}
\def\Sch{\mathit{Sch}}
\def\Mor{\mathop{\mathrm{Mor}}\nolimits}
\def\Ob{\mathop{\mathrm{Ob}}\nolimits}
\def\Sh{\mathop{\mathit{Sh}}\nolimits}
\def\NL{\mathop{N\!L}\nolimits}
\def\CH{\mathop{\mathrm{CH}}\nolimits}
\def\proetale{{pro\text{-}\acute{e}tale}}
\def\etale{{\acute{e}tale}}
\def\QCoh{\mathit{QCoh}}
\def\Ker{\mathop{\mathrm{Ker}}}
\def\Im{\mathop{\mathrm{Im}}}
\def\Coker{\mathop{\mathrm{Coker}}}
\def\Coim{\mathop{\mathrm{Coim}}}

% Boxtimes
%
\DeclareMathSymbol{\boxtimes}{\mathbin}{AMSa}{"02}

%
% Macros for moduli stacks/spaces
%
\def\QCohstack{\mathcal{QC}\!\mathit{oh}}
\def\Cohstack{\mathcal{C}\!\mathit{oh}}
\def\Spacesstack{\mathcal{S}\!\mathit{paces}}
\def\Quotfunctor{\mathrm{Quot}}
\def\Hilbfunctor{\mathrm{Hilb}}
\def\Curvesstack{\mathcal{C}\!\mathit{urves}}
\def\Polarizedstack{\mathcal{P}\!\mathit{olarized}}
\def\Complexesstack{\mathcal{C}\!\mathit{omplexes}}
% \Pic is the operator that assigns to X its picard group, usage \Pic(X)
% \Picardstack_{X/B} denotes the Picard stack of X over B
% \Picardfunctor_{X/B} denotes the Picard functor of X over B
\def\Pic{\mathop{\mathrm{Pic}}\nolimits}
\def\Picardstack{\mathcal{P}\!\mathit{ic}}
\def\Picardfunctor{\mathrm{Pic}}
\def\Deformationcategory{\mathcal{D}\!\mathit{ef}}
\usepackage{cleveref}
\crefname{originalreference}{original reference}{original references}
\Crefname{originalreference}{Original reference}{Original references}
\makeatletter
\newlabel{chow-lemma-additivity-sheaf-cycle}{{02QZ}{source}{}{sourcepointer-1}{}}
\newlabel{chow-lemma-additivity-sheaf-cycle@cref}{{[originalreference][0][]02QZ}{[1][0][]source}{}{}{}}
\makeatother
\crefname{originalreference}{original reference}{original references}
\Crefname{originalreference}{Original reference}{Original references}
\makeatletter
\newlabel{algebra-lemma-filter-Noetherian-module}{{00L0}{source}{}{sourcepointer-2}{}}
\newlabel{algebra-lemma-filter-Noetherian-module@cref}{{[originalreference][0][]00L0}{[1][0][]source}{}{}{}}
\makeatother
\crefname{originallemma}{original lemma}{original lemmas}
\Crefname{originallemma}{Original lemma}{Original lemmas}
\makeatletter
\newlabel{lemma-tor-and-diagonal}{{0B0T}{source}{}{sourcepointer-3}{}}
\newlabel{lemma-tor-and-diagonal@cref}{{[originallemma][0][]0B0T}{[1][0][]source}{}{}{}}
\makeatother
\crefname{originallemma}{original lemma}{original lemmas}
\Crefname{originallemma}{Original lemma}{Original lemmas}
\makeatletter
\newlabel{lemma-diagonal-regular-immersion}{{0B0Q}{source}{}{sourcepointer-4}{}}
\newlabel{lemma-diagonal-regular-immersion@cref}{{[originallemma][0][]0B0Q}{[1][0][]source}{}{}{}}
\makeatother
\crefname{originallemma}{original lemma}{original lemmas}
\Crefname{originallemma}{Original lemma}{Original lemmas}
\makeatletter
\newlabel{lemma-one-ideal-ci}{{0B03}{source}{}{sourcepointer-5}{}}
\newlabel{lemma-one-ideal-ci@cref}{{[originallemma][0][]0B03}{[1][0][]source}{}{}{}}
\makeatother
\crefname{originalreference}{original reference}{original references}
\Crefname{originalreference}{Original reference}{Original references}
\makeatletter
\newlabel{algebra-proposition-dimension}{{00KQ}{source}{}{sourcepointer-6}{}}
\newlabel{algebra-proposition-dimension@cref}{{[originalreference][0][]00KQ}{[1][0][]source}{}{}{}}
\makeatother
\title{Serre formula: intersection products of sheaf cycles via alternating Tor sums}
\author{The Stacks Project Authors}
\date{stacks-project:intersection.tex:local-snapshot-20261012; GNU FDL 1.2 or later}
\begin{document}
\maketitle
\noindent\textit{Editorial note.} All delivered mathematical text is the author's own, copied verbatim from the pinned Stacks Project chapter file intersection.tex; no statement or proof step was rewritten, reordered or removed.\par\medskip
\noindent\textit{Editorial note.} Span recomputation by label pairing (LANE\_BRIEF 10.2): in this chapter the author's proof environment follows the statement environment as a standalone block; the spec covers the complete environment pair (statement span, proof span) as recorded in the pool candidate.\par\medskip
\noindent\textit{Editorial note.} Occupation scan: combined claim\_key + statement\_anchor + generic-line scan over all registration files (active and withdrawn) found 0 hits for this chapter (is-occupied.json, 2026-10-14); consistent with the duty manager's six-check dispatch.\par\medskip
\ifdefined\hypertarget\hypertarget{sourcepointer-1}{}\fi
\noindent\textit{Source reference.} Original reference 02QZ: sources/stacks-chow/tags.txt, source lines 3606--3606. Cross-chapter standard prerequisite cited by the author inside the counted span; it is used as a pointer and its own proof is not part of this item.\par\medskip
\ifdefined\hypertarget\hypertarget{sourcepointer-2}{}\fi
\noindent\textit{Source reference.} Original reference 00L0: sources/stacks-chow/tags.txt, source lines 770--770. Cross-chapter standard prerequisite cited by the author inside the counted span; it is used as a pointer and its own proof is not part of this item.\par\medskip
\ifdefined\hypertarget\hypertarget{sourcepointer-3}{}\fi
\noindent\textit{Source reference.} Original lemma 0B0T: sources/stacks-chow/tags.txt, source lines 13871--13871. Cross-chapter standard prerequisite cited by the author inside the counted span; it is used as a pointer and its own proof is not part of this item.\par\medskip
\ifdefined\hypertarget\hypertarget{sourcepointer-4}{}\fi
\noindent\textit{Source reference.} Original lemma 0B0Q: sources/stacks-chow/tags.txt, source lines 13868--13868. Cross-chapter standard prerequisite cited by the author inside the counted span; it is used as a pointer and its own proof is not part of this item.\par\medskip
\ifdefined\hypertarget\hypertarget{sourcepointer-5}{}\fi
\noindent\textit{Source reference.} Original lemma 0B03: sources/stacks-chow/tags.txt, source lines 13846--13846. Cross-chapter standard prerequisite cited by the author inside the counted span; it is used as a pointer and its own proof is not part of this item.\par\medskip
\ifdefined\hypertarget\hypertarget{sourcepointer-6}{}\fi
\noindent\textit{Source reference.} Original reference 00KQ: sources/stacks-chow/tags.txt, source lines 760--760. Cross-chapter standard prerequisite cited by the author inside the counted span; it is used as a pointer and its own proof is not part of this item.\par\medskip
\begin{lemma}
\label{lemma-tor-sheaf}
\begin{reference}
\cite[Chapter V]{Serre_algebre_locale}
\end{reference}
Let $X$ be a nonsingular variety. Let $\mathcal{F}$ and
$\mathcal{G}$ be coherent sheaves on $X$ with
$\dim(\text{Supp}(\mathcal{F})) \leq r$,
$\dim(\text{Supp}(\mathcal{G})) \leq s$, and
$\dim(\text{Supp}(\mathcal{F}) \cap \text{Supp}(\mathcal{G}) )
\leq r + s - \dim X$. In this case $[\mathcal{F}]_r$ and $[\mathcal{G}]_s$
intersect properly and
$$
[\mathcal{F}]_r \cdot [\mathcal{G}]_s =
\sum (-1)^p
[\text{Tor}_p^{\mathcal{O}_X}(\mathcal{F}, \mathcal{G})]_{r + s - \dim(X)}.
$$
\end{lemma}

\begin{proof}
The statement that $[\mathcal{F}]_r$ and $[\mathcal{G}]_s$ intersect properly
is immediate. Since we are proving an equality of cycles we may work
locally on $X$. (Observe that the formation of the intersection
product of cycles, the formation of $\text{Tor}$-sheaves, and
forming the cycle associated to a coherent sheaf, each commute with
restriction to open subschemes.) Thus we may and do assume that $X$ is affine.

\medskip\noindent
Denote
$$
RHS(\mathcal{F}, \mathcal{G}) = [\mathcal{F}]_r \cdot [\mathcal{G}]_s
\quad\text{and}\quad
LHS(\mathcal{F}, \mathcal{G}) = \sum (-1)^p
[\text{Tor}_p^{\mathcal{O}_X}(\mathcal{F}, \mathcal{G})]_{r + s - \dim(X)}
$$
Consider a short exact sequence
$$
0 \to \mathcal{F}_1 \to \mathcal{F}_2 \to \mathcal{F}_3 \to 0
$$
of coherent sheaves on $X$ with
$\text{Supp}(\mathcal{F}_i) \subset \text{Supp}(\mathcal{F})$,
then both $LHS(\mathcal{F}_i, \mathcal{G})$ and
$RHS(\mathcal{F}_i, \mathcal{G})$ are defined for $i = 1, 2, 3$
and we have
$$
RHS(\mathcal{F}_2, \mathcal{G}) =
RHS(\mathcal{F}_1, \mathcal{G}) + RHS(\mathcal{F}_3, \mathcal{G})
$$
and similarly for LHS. Namely, the support condition guarantees that
everything is defined, the short exact sequence and additivity of lengths
gives
$$
[\mathcal{F}_2]_r = [\mathcal{F}_1]_r  + [\mathcal{F}_3]_r
$$
(Chow Homology, Lemma \ref{chow-lemma-additivity-sheaf-cycle})
which implies additivity for RHS. The long exact sequence of $\text{Tor}$s
$$
\ldots \to \text{Tor}_1(\mathcal{F}_3, \mathcal{G}) \to
\text{Tor}_0(\mathcal{F}_1, \mathcal{G}) \to
\text{Tor}_0(\mathcal{F}_2, \mathcal{G}) \to
\text{Tor}_0(\mathcal{F}_3, \mathcal{G}) \to 0
$$
and additivity of lengths as before implies additivity for LHS.

\medskip\noindent
By Algebra, Lemma \ref{algebra-lemma-filter-Noetherian-module}
and the fact that $X$ is affine, we can find a filtration of $\mathcal{F}$
whose graded pieces are structure sheaves of closed subvarieties of
$\text{Supp}(\mathcal{F})$. The additivity shown in the previous paragraph,
implies that it suffices to prove $LHS = RHS$ with
$\mathcal{F}$ replaced by $\mathcal{O}_V$ where
$V \subset \text{Supp}(\mathcal{F})$.
By symmetry we can do the same for $\mathcal{G}$.
This reduces us to proving that
$$
LHS(\mathcal{O}_V, \mathcal{O}_W) = RHS(\mathcal{O}_V, \mathcal{O}_W)
$$
where $W \subset \text{Supp}(\mathcal{G})$ is a closed subvariety.
If $\dim(V) = r$ and $\dim(W) = s$, then this equality is the
{\bf definition} of $V \cdot W$. On the other hand, if
$\dim(V) < r$ or $\dim(W) < s$, i.e., $[V]_r = 0$ or $[W]_s = 0$,
then we have to prove that $RHS(\mathcal{O}_V, \mathcal{O}_W) = 0$
\footnote{The reader can see that this is not a triviality by
taking $r = s = 1$ and $X$ a nonsingular surface and $V = W$
a closed point $x$ of $X$. In this case there are $3$ nonzero
$\text{Tor}$s of lengths $1, 2, 1$ at $x$.}.

\medskip\noindent
Let $Z \subset V \cap W$ be an irreducible component of dimension
$r + s - \dim(X)$. This is the maximal dimension of a component
and it suffices to show that the coefficient of $Z$ in $RHS$ is zero.
Let $\xi \in Z$ be the generic point. Write $A = \mathcal{O}_{X, \xi}$,
$B = \mathcal{O}_{X \times X, \Delta(\xi)}$, and
$C = \mathcal{O}_{V \times W, \Delta(\xi)}$.
By Lemma \ref{lemma-tor-and-diagonal} we have
$$
\text{coeff of }Z\text{ in }
RHS(\mathcal{O}_V, \mathcal{O}_W) = 
\sum (-1)^i
\text{length}_B \text{Tor}_i^B(A, C)
$$
Since $\dim(V) < r$ or $\dim(W) < s$ we have $\dim(V \times W) < r + s$
which implies $\dim(C) < \dim(X)$ (small detail omitted). Moreover, the
kernel $I$ of $B \to A$ is generated by a regular sequence of
length $\dim(X)$ (Lemma \ref{lemma-diagonal-regular-immersion}).
Hence vanishing by Lemma \ref{lemma-one-ideal-ci} because
the Hilbert function of $C$ with respect to $I$ has degree $\dim(C) < n$
by Algebra, Proposition \ref{algebra-proposition-dimension}.
\end{proof}

\end{document}
